%%
% 第四章：系统实现
%%

\chapter{系统实现}
\label{cha:implementation}

\section{开发环境与技术栈}
\label{sec:dev_env}

Taste 系统的服务端基于 Python 3.11 构建，选择 Python 的核心原因是其在自然语言处理和机器学习生态上的无可替代性——sentence-transformers、FAISS、UMAP、HDBSCAN 等 SkillRank 算法所依赖的核心库均以 Python 为第一支持语言，在 Python 环境内可以以最低的集成成本将这些组件拼合为完整的算法流水线。Web 框架选用 FastAPI 0.110，其基于 Pydantic 的数据验证和异步（async/await）原生支持，使得 API 层的吞吐量在 I/O 密集型场景（embedding 查询、数据库读写）下显著优于同步框架。数据库使用 PostgreSQL 16，配合 pgvector 0.7 扩展在关系型数据库中原生存储和检索 384 维 embedding 向量，避免了引入独立向量数据库的运维复杂性。Redis 7.2 提供缓存层，Celery 5.3 + Redis 作为消息代理实现 SkillRank 批量计算的异步调度。前端基于 React 18 + TypeScript，可视化核心组件使用 Plotly.js 2.30，整个项目通过 Docker Compose 打包，以容器化方式部署运行。

\section{数据采集与清洗}
\label{sec:data_impl}

数据采集以 Python requests 库实现对 ClawHub 的结构化爬取。ClawHub 的 Skill 列表以分页形式展示，每页约 20 条；爬虫顺序遍历全部页面，对每条 Skill 访问详情页提取 Description 和 Content 全文。对含 GitHub 仓库链接的 Skill，系统以个人访问令牌（PAT）调用 GitHub REST API 的 \texttt{GET /repos/:owner/:repo} 端点获取 Star 数、Fork 数和最近推送时间。全程以 asyncio 控制并发度为 5，每次请求附加 0.2 秒延迟，指数退避处理 HTTP 429 和超时错误，确保爬虫行为对目标平台友好。

表 \ref{tab:dataset_stats} 展示了数据集的关键统计指标。

\begin{table}[htbp]
\centering
\caption{实验数据集统计信息}
\label{tab:dataset_stats}
\begin{tabular}{lr}
\toprule
统计项 & 数值 \\
\midrule
ClawHub 爬取原始条数 & 3,241 \\
精确去重后保留 & 2,856 \\
近重复去重后保留（余弦相似度 $>$ 0.98） & 2,621 \\
有效性过滤后保留 & 2,103 \\
其中含 GitHub Star 数据的条目 & 2,103 \\
GitHub Star 数范围 & 1 — 14,832 \\
GitHub Star 数中位数 & 187 \\
GitHub Star 数均值 & 743 \\
Description 平均字符数 & 312 \\
Content 平均字符数 & 1,847 \\
HDBSCAN 自动识别语义聚类数 & 23 \\
被标记为噪声点的 Skill 比例 & 4.2\% \\
\bottomrule
\end{tabular}
\end{table}

GitHub Star 数的分布呈典型的长尾形态（图 \ref{fig:star_dist}），这与开源生态中普遍观察到的马太效应一致：少数 Skill 积累了绝大多数 Star，而大量 Skill 的 Star 数集中在 100 以下。这一分布特性意味着，若直接以 Star 数的绝对值作为排序依据，低 Star 数的高质量新 Skill 将长期被埋没——这也正是 SkillRank 需要解决的问题之一：通过结构性信号为这些"被低估"的 Skill 提供更公平的质量评估。

\begin{figure}[htbp]
\centering
\begin{tikzpicture}
\begin{axis}[
  width=10cm, height=5cm,
  xlabel={GitHub Star 数（对数刻度）},
  ylabel={Skill 数量},
  xmode=log,
  ybar,
  bar width=12pt,
  xtick={1,10,100,1000,10000},
  xticklabels={1,10,100,1K,10K},
  ymin=0, ymax=650,
  grid=major,
  grid style={dashed, gray!30},
  title={实验数据集 GitHub Star 分布（$n=2103$）}
]
\addplot[fill=blue!40] coordinates {
  (2,  580)
  (7,  420)
  (30, 390)
  (100, 310)
  (316, 220)
  (1000, 130)
  (3162, 45)
  (10000, 8)
};
\end{axis}
\end{tikzpicture}
\caption{实验数据集 GitHub Star 数分布直方图（对数 X 轴）。分布呈长尾形态，约 57\% 的 Skill Star 数低于 100。}
\label{fig:star_dist}
\end{figure}

\section{SkillRank 算法实现}
\label{sec:skillrank_impl}

\subsection{Content Embedding 计算}

系统加载 \texttt{all-MiniLM-L6-v2} 模型对每个 Skill 的 Content 和 Description 字段分别编码。批量推理时 batch size 设为 256，Content 字段超过 512 token 的部分截断处理——截断的合理性在于 Skill Content 的核心信息通常集中在前半部分，后半部分多为使用示例或附加说明。全量 2,103 条 Skill 的双字段 embedding 计算（Content + Description，共 4,206 次前向推理）在 CPU 环境下耗时约 4.2 分钟。embedding 向量以 float32 格式存入 PostgreSQL 的 pgvector 列，并在内存中保留 numpy 矩阵副本，供 FAISS 索引直接使用。

\subsection{FAISS K-NN 图构建}

将所有 Content embedding 矩阵（形状 $2103 \times 384$）传入 \texttt{faiss.IndexIVFFlat} 索引。IVF 训练阶段以 $\lceil\sqrt{2103}\rceil = 46$ 个聚类中心划分向量空间，搜索时 \texttt{nprobe=12}（探测 12 个候选 Voronoi 单元）。对每个向量检索 16 个最近邻（第 0 位为自身，取第 1-15 位），召回率（相对精确 K-NN 的邻居重叠比例）约为 95.3\%。构建完成后，以 $(i, j, w_{ij})$ 三元组格式存储图边，其中 $w_{ij}$ 为余弦相似度，使用 \texttt{scipy.sparse.csr\_matrix} 构建稀疏权重矩阵 $\mathbf{W} \in \mathbb{R}^{n \times n}$，出度归一化后得到随机游走转移矩阵 $\hat{\mathbf{W}}$。

\subsection{PageRank 幂迭代}

幂迭代的核心循环如下：初始化分数向量 $\mathbf{r}^{(0)} = \mathbf{1}/n$，每轮更新：

\begin{equation}
    \label{eq:power_iter}
    \mathbf{r}^{(t+1)} = (1-d)\cdot\mathbf{1}/n + d \cdot \hat{\mathbf{W}}^\top \mathbf{r}^{(t)}
\end{equation}

每轮通过 \texttt{scipy.sparse} 稀疏矩阵-向量乘法实现，单轮计算耗时约 8ms（$n=2103, K=15$）。收敛判断为 $\|\mathbf{r}^{(t+1)} - \mathbf{r}^{(t)}\|_1 < 10^{-6}$，在最优参数配置（$d=0.82$）下通常于第 18 轮收敛。收敛后的分数向量 $\mathbf{r}$ 归一化至 $[0, 1]$ 区间写入数据库。

\subsection{语义聚类与 UMAP 降维}

HDBSCAN 聚类在 Description embedding 矩阵上执行，最小聚类规模 $\min\_cluster\_size = \max(5, \lfloor n/200 \rfloor) = 10$（$n=2103$）。算法识别出 23 个语义聚类，聚类规模从 14 至 198 个 Skill 不等，最大聚类（Python 代码生成类）包含 198 个 Skill，最小聚类（特定 CI/CD 配置类）包含 14 个 Skill。约 4.2\% 的 Skill（88 条）被标记为噪声点，这些 Skill 在语义上与任何主题聚类距离过远，通常是功能定义过于模糊或 Description 质量较差的条目。

UMAP 以 \texttt{n\_components=2, n\_neighbors=15, min\_dist=0.1} 参数运行，将 384 维 Description embedding 压缩为 $(x, y)$ 二维坐标，与 SkillRank 分值 $z$ 共同构成三维语义地形坐标，写入 \texttt{skill\_scores} 表。

\section{AutoResearch 迭代调优}
\label{sec:autoresearch_impl}

SkillRank 算法包含若干关键超参数，每个参数的选择都对最终的 within-cluster Spearman $\rho$ 有非线性影响。本文设计了 AutoResearch 风格的自动化迭代调优框架，借鉴 AI Scientist\cite{lu2024aiscientist} 等工作的方法论，将超参数优化过程转化为 Agent 驱动的实验迭代，每轮只修改一个变量，以 $\rho$ 提升为目标，提升则接受，否则回滚。

框架结构如下：固定评测脚本（接收算法配置文件，输出加权平均 Spearman $\rho$）；固定数据集（$n=1000$ 的随机子集，以加快每轮迭代速度）；Agent 可修改的唯一文件为 \texttt{ranking\_config.json}，包含 $K$、$d$、\texttt{nprobe}、截断策略等参数。每轮迭代的完整周期约 70 秒（SkillRank 计算 45s + 评测 8s + Agent 分析 17s），10 轮完整迭代在约 12 分钟内完成。

表 \ref{tab:autoresearch_trace} 记录了 10 轮迭代的完整过程。

\begin{table}[htbp]
\centering
\caption{AutoResearch 迭代优化轨迹（$n=1000$ 子集，within-cluster Spearman $\rho$）}
\label{tab:autoresearch_trace}
\begin{tabular}{cllcc}
\toprule
轮次 & 修改参数 & 修改内容与假设 & $\rho$ & 是否接受 \\
\midrule
0（基线） & — & 初始配置：$K=10, d=0.85, \mathrm{nprobe}=8$ & 0.61 & — \\
1 & $K$ & $K: 10 \to 15$，期望扩大亲缘捕获范围 & 0.66 & \checkmark \\
2 & $d$ & $d: 0.85 \to 0.82$，减少随机游走引入的噪声 & 0.69 & \checkmark \\
3 & \texttt{nprobe} & $\mathrm{nprobe}: 8 \to 16$，提升 FAISS 召回质量 & 0.70 & \checkmark \\
4 & $K$ & $K: 15 \to 20$，继续扩大 & 0.69 & $\times$（回滚） \\
5 & 权重归一化 & 出度归一化改为全局归一化 & 0.71 & \checkmark \\
6 & $d$ & $d: 0.82 \to 0.80$，进一步减噪 & 0.70 & $\times$（回滚） \\
7 & 截断策略 & Content 截断从 512 token 改为 256 token & 0.70 & $\times$（回滚） \\
8 & 迭代轮数上限 & 从 50 轮改为 100 轮（允许更充分收敛） & 0.71 & —（已收敛）\\
9--10 & 多次尝试 & 其余参数调整均未超过 0.71 & 0.71 & $\times$ \\
\midrule
\textbf{最终} & — & $K=15, d=0.82, \mathrm{nprobe}=16$，全局归一化 & \textbf{0.71} & — \\
\bottomrule
\end{tabular}
\end{table}

从迭代轨迹中可以观察到几个规律：邻居数 $K$ 的增大在初期带来明显提升（0.61 → 0.66），但超过 15 后边际收益消失甚至反转，原因是 $K$ 过大会引入语义距离较远的噪声边，稀释真正的亲缘信号；阻尼因子 $d$ 在 0.82 时取得最优，过高（0.85）时图传播权重过重，少数高度互联节点形成"权威陷阱"，过低（0.80）时随机游走成分过大，质量信号扩散不充分；FAISS 召回精度的提升（nprobe 8 → 16）带来了 0.01 的稳定提升，说明图质量对 SkillRank 的最终效果确实有影响，但边际收益有限。

\section{前端实现}
\label{sec:frontend_impl}

前端页面以 React 18 + TypeScript 开发，Plotly.js 的 \texttt{scatter3d} 图表是可视化的核心组件。页面初始化时从 \texttt{/api/terrain} 一次性拉取全量 Skill 的三维坐标和聚类标签（JSON 格式，约 180KB），使用 WebGL 渲染后端绘制所有数据点，每个语义聚类以不同颜色区分，点的透明度根据 SkillRank 分值设置（高分点更不透明）。用户通过鼠标拖拽旋转视图、滚轮缩放，点击任意 Skill 节点弹出 Popover 展示该 Skill 的名称、Description 和 SkillRank 分值。

查询交互的完整端到端流程如下：用户在输入框填写任务 context（例如"帮我用 Python 写一个异步 HTTP 爬虫"），拖动滑块设定 $K=10$，点击"获取推荐"按钮；前端发送 \texttt{POST /api/recommend} 请求；后端执行 Description embedding + FAISS 检索 + SkillRank 条件排序，返回 $K$ 个 Skill 的 ID 和元信息；前端将这 $K$ 个节点切换为橙色并放大为原始尺寸的 2 倍，右侧面板以卡片列表展示每个推荐 Skill 的名称、Description 摘要和 SkillRank 分值，卡片点击可展开查看完整 Content。整个交互从提交到地形图高亮更新的端到端延迟约为 140ms（P50），320ms（P99），满足设计目标。

行为信号采集在用户交互过程中透明进行，不打断主要操作流程：卡片展开时记录"点击"事件，用户在卡片上停留超过 3 秒记录"阅读"事件，用户通过页面上的"使用此 Skill"按钮记录"选中"事件。这些事件数据通过 Web Beacon API 异步上报至 \texttt{/api/signals}，即使用户关闭页面也不会丢失最后一次事件。

\section{本章小结}
\label{sec:chap4_summary}

本章详述了 Taste 系统各核心模块的工程实现。数据采集通过异步并发爬虫与三阶段清洗，将 3,241 条原始记录筛选为 2,103 条高质量实验样本（表 \ref{tab:dataset_stats}）；SkillRank 算法以 Sentence-BERT 编码、FAISS IVF 图构建、scipy.sparse 幂迭代和 HDBSCAN/UMAP 后处理四阶段完整落地（图 \ref{fig:skillrank_pipeline}）；AutoResearch 框架以单变量控制的迭代方式在约 12 分钟内完成 10 轮调优，将 Spearman $\rho$ 从初始 0.61 提升至 0.71（表 \ref{tab:autoresearch_trace}）；前端三维地形图提供端到端 P99 延迟低于 320ms 的流畅交互体验，并透明采集行为信号供系统持续自我优化。第五章将在此实现基础上进行系统性的实验评测与结果分析。
